%===========================
% Chapitre 2 : Classification Under Demographic Parity Constraint
%===========================
\chapter{Classification Under Demographic Parity Constraint}

\section{Notations}

Considérons les notations suivantes pour ce chapitre :
\begin{itemize}
    \item Une observation est un triplet $(X, S, Y)$.
    \item $X \in \mathbb{R}^d$ représente les caractéristiques (features).
    \item $S \in \{-1, 1\}$ est l'attribut sensible (sensitive attribute).
    \item $Y \in \{0, 1\}$ est la sortie (label).
\end{itemize}

Un prédicteur (ou classifieur) $f$ est une fonction de $(X,S)$ telle que $f(X,S) \in \{0, 1\}$.

Nous considérons le risque de mauvaise classification (\textit{misclassification risk}) :
\begin{equation}
    R(f) = \mathbb{P}(f(X,S) \neq Y)
\end{equation}
comme mesure de performance.

Dans ce cas, le prédicteur optimal (ou Classifieur de Bayes) qui satisfait :
\begin{equation}
    f^* \in \argmin_{f} R(f)
\end{equation}
est défini par :
\begin{equation}
    f^*(X,S) = \indic\{\eta(X,S) \ge 1/2\}
\end{equation}
où $\eta(X,S) = \mathbb{P}(Y=1 | X,S)$.

\section{Définition de l'Équité (Fairness)}

Nous considérons la contrainte de Parité Démographique (Demographic Parity constraint ou DP-constraint).
Nous disons qu'un prédicteur $f$ satisfait la contrainte DP si :
\begin{equation}
    \forall s \in \{-1, 1\}, \quad \mathbb{P}(f(X,S)=1 | S=s) = \mathbb{P}(f(X,S)=1)
\end{equation}
Ceci équivaut aux propriétés suivantes :
\begin{align}
    f(X,S) \perp \!\!\! \perp S \quad &\iff \mathbb{P}(f(X,S)=1 | S=1) = \mathbb{P}(f(X,S)=1 | S=-1) \\
    &\iff \sum_{s \in \{-1, 1\}} s \cdot \mathbb{P}(f(X,S)=1 | S=s) = 0
\end{align}

\section{Classifieur Optimal sous Contrainte DP}

Nous considérons $f^*_{fair}$ défini comme :
\begin{equation}
    f^*_{fair} \in \argmin_{f \text{ satisfaisant DP}} R(f)
\end{equation}
Ceci est équivalent à résoudre le problème d'optimisation sous contrainte :
\begin{equation}
    f^*_{fair} \in \argmin_{\substack{f \\ \text{t.q. } \sum_{s \in \{-1, 1\}} s \cdot \mathbb{P}(f(X,S)=1 | S=s) = 0}} R(f)
\end{equation}

Le Lagrangien associé s'écrit comme suit :
\begin{equation}
    \mathcal{L}(f, \lambda) = R(f) + \lambda \sum_{s \in \{-1, 1\}} s \mathbb{P}(f(X,S)=1 | S=s)
\end{equation}
En introduisant la notation $\mathbb{P}(\cdot | S=s) := \mathbb{P}_{X|S=s}(\cdot)$, on a :
\begin{equation}
    \mathcal{L}(f, \lambda) = R(f) + \lambda \sum_{s \in \{-1, 1\}} s \mathbb{P}_{X|S=s}(f(X,S)=1)
\end{equation}

Pour résoudre le problème d'optimisation :
\begin{equation}
    f^*_{fair} \in \argmin_{f \text{ satisfaisant DP}} R(f)
\end{equation}
nous devons résoudre le problème d'optimisation suivant (problème primal) :
\begin{equation}
    \min_{f} \max_{\lambda} \mathcal{L}(f, \lambda)
\end{equation}

En général, nous résolvons le problème dual :
\begin{equation} \label{eq:dual_problem} \tag{$*$}
    \max_{\lambda} \min_{f} \mathcal{L}(f, \lambda)
\end{equation}
Nous avons toujours la propriété suivante (\textit{principe de dualité faible}) :
\begin{equation}
    \min_{f} \max_{\lambda} \mathcal{L}(f, \lambda) \ge \max_{\lambda} \min_{f} \mathcal{L}(f, \lambda)
\end{equation}
Nous résolvons le problème \eqref{eq:dual_problem} et montrons que :
\begin{equation}
    \min_{f} \max_{\lambda} \mathcal{L}(f, \lambda) = \max_{\lambda} \min_{f} \mathcal{L}(f, \lambda)
\end{equation}
C'est-à-dire qu'il n'y a pas de saut de dualité (\textit{no duality gap}, ou dualité forte).

\section{Résolution du Problème Dual}

Le but est d'obtenir une expression analytique (\textit{closed form expression}) du prédicteur optimal $f^*_{fair}$.

\subsection*{Première étape : Minimisation par rapport à $f$}

Pour tout $\lambda \in \mathbb{R}$ fixé, nous cherchons à résoudre :
\begin{equation}
    \min_{f} \mathcal{L}(f, \lambda)
\end{equation}
Rappelons l'expression du Lagrangien :
\begin{equation}
    \mathcal{L}(f, \lambda) = R(f) + \lambda \sum_{s \in \{-1, 1\}} s \mathbb{P}_{X|S=s}(f(X,S)=1)
\end{equation}
Le risque s'écrit comme une espérance :
\begin{align}
    R(f) &= \mathbb{P}(f(X,S) \neq Y) \\
    &= \mathbb{E}[\indic\{f(X,S) \neq Y\}]
\end{align}
En conditionnant par rapport à $Y$, on obtient :
\begin{align}
    R(f) &= \mathbb{E}\left[ \indic\{f(X,S) \neq Y\} \indic\{Y=0\} \right] + \mathbb{E}\left[ \indic\{f(X,S) \neq Y\} \indic\{Y=1\} \right]
\end{align}
Pour tout $y \in \{0, 1\}$, en utilisant la loi des espérances itérées par rapport à $(X,S)$ :
\begin{align}
    \mathbb{E}\left[ \indic\{f(X,S) \neq y\} \indic\{Y=y\} \right] &= \mathbb{E}\left[ \mathbb{E}\left[ \indic\{f(X,S) \neq y\} \indic\{Y=y\} \big| X, S \right] \right] \\
    &= \mathbb{E}\left[ \indic\{f(X,S) \neq y\} \mathbb{E}\left[ \indic\{Y=y\} \big| X, S \right] \right] \\
    &= \mathbb{E}\left[ \indic\{f(X,S) \neq y\} \mathbb{P}(Y=y | X, S) \right]
\end{align}
En notant $\eta(X,S) = \mathbb{P}(Y=1 | X,S)$, on a $\mathbb{P}(Y=0 | X,S) = 1 - \eta(X,S)$.
Ainsi, le risque devient :
\begin{align}
    R(f) &= \mathbb{E}\left[ \indic\{f(X,S) \neq 0\} (1 - \eta(X,S)) \right] + \mathbb{E}\left[ \indic\{f(X,S) \neq 1\} \eta(X,S) \right]
\end{align}
Puisque $f(X,S) \in \{0, 1\}$, nous avons les identités suivantes :
\begin{align}
    \indic\{f(X,S) \neq 0\} &= \indic\{f(X,S) = 1\} = f(X,S) \\
    \indic\{f(X,S) \neq 1\} &= \indic\{f(X,S) = 0\} = 1 - f(X,S)
\end{align}
En substituant ces identités dans l'expression du risque, on obtient :
\begin{align}
    R(f) &= \mathbb{E}\left[ f(X,S) (1 - \eta(X,S)) \right] + \mathbb{E}\left[ (1 - f(X,S)) \eta(X,S) \right] \\
    &= \mathbb{E}\left[ f(X,S) (1 - 2\eta(X,S)) \right] + \mathbb{E}\left[ \eta(X,S) \right]
\end{align}
